% ======================================================================
%  sec25.tex —— 第 25 节 Proj(S)
%  底本：原书 pp.119–123；原书标题：25. Proj(S)
% ======================================================================

\section{$\Proj(S)$}\label{sec:25}

\subsection{$\Proj(S)$ 的底空间}\label{sec:25.1}

本小节固定一个 $\Z_{\ge 0}$-分次（交换）环 $S$。换言之，$S$ 是一个交换环，作为阿贝尔群带有直和分解
\[
S = \bigoplus_{n \in \Z_{\ge 0}} S_n,
\]
并且满足：
\begin{itemize}[leftmargin=2em,itemsep=0.25ex]
  \item 单位元 $1$ 的次数为 $0$，即 $1 \in S_0$；
  \item 乘法与分次相容，即对 $x \in S_m$、$y \in S_n$ 有 $xy \in S_{m+n}$。
\end{itemize}

回忆：元素 $x \in S$ 称为\emph{齐次}的，若对某个 $n \in \Z$ 有 $x \in S_n$。理想 $I \subseteq S$ 称为\emph{齐次}的，若它可由齐次元素生成；等价地，$I$ 齐次当且仅当 $I \simeq \bigoplus_{n \in \Z_{\ge 0}} (I \cap S_n)$。

\begin{example}\label{ex:25.1.1}
\emph{增广理想}
\[
S_+ := \bigoplus_{n>0} S_n = \ker(S \to S_0)
\]
是齐次理想。
\end{example}

\begin{definition}\label{def:25.1.2}
我们定义 $S$ 的\emph{齐次（素）谱}为如下子空间
\[
\Proj(S) \subseteq \Spec(S),
\]
它由那些对应于齐次素理想 $\mathfrak p \subseteq S$ 且满足 $S_+ \not\subseteq \mathfrak p$ 的点组成。
\end{definition}

\begin{yizhu}
$S_+$ 通常称为\emph{无关理想}（irrelevant ideal），$\Proj(S)$ 的定义就是「齐次素谱去掉 $S_+$ 的零点集」。条件 $S_+ \not\subseteq \mathfrak p$ 等价于存在正次数的齐次元素 $f \in S_+ \setminus \mathfrak p$，而这样的 $f$ 正是下文标准开集 $U_+(f)$ 的指标：没有它，$\Proj(S)$ 上的结构层无从定义。Stacks Project 用同一条件定义齐次谱。
[来源：Stacks Project, Tag 00JN（Definition 10.57.1，齐次谱 $\Proj(S)$ 的定义），\url{https://stacks.math.columbia.edu/tag/00JN}（2026-11-17 访问）。]
\end{yizhu}

\begin{exercise}\label{exc:25.1.3}
设 $k$ 为代数闭域，$n \ge 0$ 为整数。设 $S = k[x_1, \dots, x_n]$ 是 $\Z_{\ge 0}$-分次环，其中 $S_0 = k$、$x_i \in S_1$。证明：对任意点 $y \in \Spec(S) = \A^n_k$，下述条件等价：
\begin{enumerate}[label=(\arabic*),leftmargin=2.4em,itemsep=0.25ex]
  \item 点 $y$ 含于 $\Proj(S)$。
  \item 点 $y$ 非闭，且 $\overline{\{y\}} \cap \A^n_k(k)$ 是 $\A^n_k(k) \simeq k^{\oplus n}$ 的锥形子集。
\end{enumerate}
\end{exercise}

\begin{notation}\label{not:25.1.4}
我们使用下述记号：
\begin{itemize}[leftmargin=2em,itemsep=0.25ex]
  \item 对元素 $f \in S$，定义
  \[
  U_+(f) := \Proj(S) \cap U(f),
  \]
  它是 $\Proj(S)$ 的开子集。
  \item 对理想 $I \subseteq S$，定义
  \[
  Z_+(I) := \Proj(S) \cap Z(I),
  \]
  它是 $\Proj(S)$ 的闭子集。
\end{itemize}
\end{notation}

展开定义即得下述结果。

\begin{lemma}\label{lem:25.1.5}
设 $S$ 是 $\Z_{\ge 0}$-分次环。
\begin{itemize}[leftmargin=2em,itemsep=0.25ex]
  \item 对任意元素 $f = \sum_{i=0}^n f_i \in S$（其中 $f_i \in S_i$），有
  \[
  U_+(f) = \bigcup_{i=0}^n U_+(f_i).
  \]
  \item 对任意元素 $f_0 \in S_0$，有
  \[
  U_+(f_0) = \bigcup_{g \in S_n,\ n>0} U_+(f_0 g).
  \]
  \item 对任意理想 $I \subseteq S$，有
  \[
  Z_+(I) = Z_+\Bigl(\bigoplus_{n \in \Z_{\ge 0}} (I \cap S_n)\Bigr).
  \]
  \item 对任意理想 $I \subseteq S$ 与 $n \ge 0$，有
  \[
  Z_+(I) = Z_+\bigl(\sqrt{I} \cap S_+\bigr).
  \]
  此外，若 $I$ 齐次，则 $\sqrt{I} \cap S_+$ 也齐次。
\end{itemize}
\end{lemma}

与仿射情形一样，我们有下述结果。

\begin{corollary}\label{cor:25.1.6}
开子集 $U_+(f)$ 构成 $\Proj(S)$ 拓扑的一组基，其中 $f$ 取遍 $S_+$ 中全部齐次元素。我们称它们为 $\Proj(S)$ 的\emph{标准开子集}。
\end{corollary}

\begin{yizhu}
标准开集并不自动给出拟紧性。Stacks Project 记 $U_d = \bigcup_{f \in S_d} D_+(f) \subseteq \Proj(S)$，指出 $\Proj(S) = \bigcup_{d \ge 1} U_d$，但对一般的 $S$，$U_d$ 与 $\Proj(S)$ 本身都未必拟紧。故拟紧性必须另加条件，这正是引理 \ref{lem:25.4.2} 的由来。
同一记号下还引入 Veronese 子环 $S^{(d)}$：换成 $S^{(d)}$ 后 $\cO(d)$ 才成为线丛，这正是习题 \ref{exc:25.2.10} 的用处。

[来源：Stacks Project, Tag 01N4（Section 27.12，Morphisms into Proj；其中有 $U_d$ 的定义与「Neither $X$ nor $U_d$ need be quasi-compact」），\url{https://stacks.math.columbia.edu/tag/01N4}（2026-11-17 访问）。]
\end{yizhu}

\begin{corollary}\label{cor:25.1.7}
$\Proj(S)$ 的闭子集与满足 $I = \sqrt{I} \cap S_+$ 的齐次理想 $I \subseteq S_+$ 一一对应。
\end{corollary}

\begin{notation}\label{not:25.1.8}
设 $f \in S_+$ 为齐次元素。环 $S_f$ 带有自然的 $\Z$-分次。记 $S_{(f)} \subseteq S_f$ 为其中分次为零的元素构成的子环。
\end{notation}

\begin{remark}\label{rem:25.1.9}
对 $f \in S_n$，有
\[
S_{(f)} \simeq \Bigl(\bigoplus_{d \in \Z_{\ge 0}} S_{dn}\Bigr)/(f-1).
\]
\end{remark}

\begin{warning}\label{warn:25.1.10}
环 $S_{(f)}$ 不是 $S$ 对乘集 $S \setminus (f)$ 的局部化。事实上，$(f)$ 一般不是素理想。
\end{warning}

\begin{exercise}\label{exc:25.1.11}
证明合成映射
\[
U_+(f) \subseteq U(f) \simeq \Spec(S_f) \to \Spec(S_{(f)})
\]
是同胚。
\end{exercise}

\begin{yizhu}
同胚的实质是：$U_+(f)$ 上可用的函数形如 $g/f^d$ 且总次数为零，故坐标环取 $S_f$ 的零次部分 $S_{(f)}$，而非整个 $S_f$。Stacks Project 把标准开集记作 $D_+(f)$，证明 $D_+(f) \cong \Spec(S_{(f)})$ 并由之得到 $\Proj(S)$ 的概形结构；这与下文构造 \ref{con:25.2.2} 用 $S_f$ 的零次部分定义结构层是一回事。
[来源：Stacks Project, Tag 01MB（Lemma 27.8.7：$\Proj(S)$ 是概形，$D_+(f)$ 仿射，$\widetilde{M}$ 拟凝聚），\url{https://stacks.math.columbia.edu/tag/01MB}（2026-11-17 访问）。]
\end{yizhu}

\subsection{概形 $\Proj(S)$}\label{sec:25.2}

本小节固定一个 $\Z_{\ge 0}$-分次环 $S$。记
\[
\Spec(S)^{\circ} := \Spec(S) \setminus \Spec(S_0) = \bigcup_{f \in \bigoplus_{n>0} S_n} U(f),
\]
其中我们把 $\Spec(S_0)$ 经由投影 $S \to S_0$ 视为 $\Spec(S)$ 的闭子概形。注意 $\Spec(S)^{\circ}$ 的点正是那些不含 $S_+$ 的（可能非齐次的）素理想 $\mathfrak p \subseteq S$。

\begin{exercise}\label{exc:25.2.1}
映射
\[
\pi \colon \Spec(S)^{\circ} \to \Proj(S), \qquad \mathfrak p \mapsto \bigoplus_{n \in \Z_{\ge 0}} (\mathfrak p \cap S_n)
\]
有定义且连续。此外，对任意齐次 $f \in S_+$，有 $\pi^{-1}(U_+(f)) = U(f)$。
\end{exercise}

\begin{construction}\label{con:25.2.2}
考虑对象
\[
\pi_* \cO_{\Spec(S)^{\circ}} \in \Shv(\Proj(S), \CRing).
\]
对任意标准开子集 $V \subseteq \Proj(S)$，取齐次元素 $f \in S_+$ 使 $V = U_+(f)$。则
\[
(\pi_* \cO_{\Spec(S)^{\circ}})(V) \simeq \cO_{\Spec(S)^{\circ}}(U(f)) \simeq S_f.
\]
\end{construction}

这在 $(\pi_* \cO_{\Spec(S)^{\circ}})(V)$ 上给出一个 $\Z$-分次。可以验证该分次不依赖 $f$ 的选取；而且对标准开子集 $V_1 \subseteq V_2$，限制映射
\[
(\pi_* \cO_{\Spec(S)^{\circ}})(V_2) \to (\pi_* \cO_{\Spec(S)^{\circ}})(V_1)
\]
是 $\Z$-分次同态。因此我们得到一个升级
\[
\pi_* \cO_{\Spec(S)^{\circ}} \in \Shv(\Proj(S), \CRing_{\mathrm{gr}}),
\]
其中 $\CRing_{\mathrm{gr}}$ 是 $\Z$-分次环的范畴。今定义
\[
\cO_{\Proj(S)} := (\pi_* \cO_{\Spec(S)^{\circ}})_0 \in \Shv(\Proj(S), \CRing)
\]
为该对象的零次部分，并称之为 $\Proj(S)$ 的\emph{结构层}。

我们滥用记号，仍用 $\Proj(S)$ 表示环化空间 $(\Proj(S), \cO_{\Proj(S)})$。注意我们有环化空间的态射 $\pi \colon \Spec(S)^{\circ} \to \Proj(S)$，它由态射
\[
\cO_{\Proj(S)} = (\pi_* \cO_{\Spec(S)^{\circ}})_0 \to \pi_* \cO_{\Spec(S)^{\circ}}
\]
给出。

\begin{remark}\label{rem:25.2.3}
由构造，在任意标准开子集 $U_+(f) \subseteq \Proj(S)$ 上，有
\[
\cO_{\Proj(S)}(U_+(f)) \simeq S_{(f)}.
\]
设 $x \in \Proj(S)$ 为一点，$\mathfrak p \subseteq S$ 为对应的齐次理想，则
\[
\cO_{\Proj(S), x} \simeq S_{(\mathfrak p)},
\]
即 $S_{\mathfrak p}$ 的零次部分。
\end{remark}

\begin{proposition}\label{prop:25.2.4}
环化空间 $(\Proj(S), \cO_{\Proj(S)})$ 是概形，且 $\pi \colon \Spec(S)^{\circ} \to \Proj(S)$ 是概形的仿射态射。
\end{proposition}

\begin{proof}
容易看出：对任意标准开子集 $U_+(f) \subseteq \Proj(S)$，环化空间 $\bigl(U_+(f),\, \cO_{\Proj(S)}|_{U_+(f)}\bigr)$ 同构于概形 $\Spec(S_{(f)})$。而且在这个开子空间上，$\pi$ 可以等同于对应环同态 $S_{(f)} \to S_f$ 的态射 $\Spec(S_f) \to \Spec(S_{(f)})$。
\end{proof}

\begin{warning}\label{warn:25.2.5}
连续映射 $\Proj(S) \to \Spec(S)$ 一般不能升级为概形态射。
\end{warning}

\begin{yizhu}
理由是仿射概形 $\Spec(S)$ 上的态射由整体截面决定：若该连续映射来自概形态射，它应当对应环同态 $S \to \Gamma(\Proj(S), \cO_{\Proj(S)})$。而由注记 \ref{rem:25.2.3}，$\Proj(S)$ 的局部环是 $S_{\mathfrak p}$ 的零次部分，正次元素在 $\Proj(S)$ 上不可见，故 $S$ 的正次部分丢失；这与警告 \ref{warn:25.2.9}（$S$ 不能由 $\Proj(S)$ 还原）互为表里。
[依据：原书 p.121 注记 25.2.3 与警告 25.2.9。]
\end{yizhu}

\begin{remark}\label{rem:25.2.6}
应当把 $\Proj(S)$ 看成 $\Spec(S)^{\circ}$ 对某个\emph{伸缩作用}的商。
\end{remark}

\begin{yizhu}
伸缩作用即乘法群 $\GG_m = \Spec \Z[t, t^{-1}]$ 的作用：$t$ 在 $S_n$ 上按 $n$ 次标量作用。经典情形见 nLab 的「tautological line bundle」条目，它把 $\PP^n$ 写成 $k^{n+1} \setminus \{0\}$ 对「相差非零纯量倍」这一等价关系的商；分次情形只是把 $k^\times$ 换成 $\GG_m$，而 $\pi$ 的每条纤维恰为一条轨道。
[来源：nLab "tautological line bundle"，\url{https://ncatlab.org/nlab/show/tautological+line+bundle}（Last revised 2026-04-06；2026-11-17 访问）。]
\end{yizhu}

\begin{exercise}\label{exc:25.2.7}
设 $\pi \colon \Spec(S)^{\circ} \to \Proj(S)$ 为典范态射。
\begin{enumerate}[label=(\arabic*),leftmargin=2.4em,itemsep=0.25ex]
  \item 证明 $\pi$ 是概形的满态射。
  \item 证明 $\Spec(S)^{\circ} \overset{\subseteq}{\longrightarrow} \Spec(S) \to \Spec(S_0)$ 唯一地穿过 $\Proj(S)$ 分解。
\end{enumerate}
\end{exercise}

特别地，若 $S$ 是 $\Z_{\ge 0}$-分次 $A$-代数，则 $\Proj(S)$ 典范地是 $A$-概形。

\begin{exercise}\label{exc:25.2.8}
设 $R$ 为交换环，$n \ge 0$。设 $S = R[x_0, \dots, x_n]$ 是 $\Z_{\ge 0}$-分次环，其中 $S_0 = R$、$x_i \in S_1$。证明 $\Proj(S) \simeq \PP^n_R$。
\end{exercise}

\begin{warning}\label{warn:25.2.9}
$\Z_{\ge 0}$-分次交换环 $S$ 不能由 $\Proj(S)$ 还原。
\end{warning}

\begin{exercise}\label{exc:25.2.10}
设 $d$ 为正整数，$S' = \bigoplus_{n \ge 0} S'_{nd}$ 是 $\Z_{\ge 0}$-分次环，其中 $S'_{nd} := S_n$。证明 $\Proj(S') \simeq \Proj(S)$。
\end{exercise}

\begin{yizhu}
取 $d$ 次 Veronese 子环 $S' = \bigoplus_n S_{nd}$ 后 $\Proj$ 不变，这与警告 \ref{warn:25.2.9} 是同一现象的两面：$\Proj(S)$ 只记得齐次元素之间的比值，故丢掉 $d$ 的倍数之外的次信息。Stacks Project 正是用 $S^{(d)}$ 与 $O(d)$（而非 $O(1)$）作相对 Proj 的一般理论，以绕开 $O(1)$ 未必可逆的困难。
[依据：原书 p.122 警告 25.2.9 与习题 25.2.10 互见；另见 Stacks Project, Tag 01N4（$S^{(d)}$ 与 $U_d$ 的定义）。]
\end{yizhu}

\subsection{函子性}\label{sec:25.3}

本小节固定 $\Z_{\ge 0}$-分次环之间的分次同态 $\phi \colon S \to S'$。

一般而言，不应期待有态射 $\Proj(S') \to \Proj(S)$，因为态射 $\Spec(S') \to \Spec(S)$ 并不把 $\Spec(S')^{\circ}$ 送进 $\Spec(S)^{\circ}$。

\begin{notation}\label{not:25.3.1}
考虑
\[
U(\phi) := \Spec(S') \times_{\Spec(S)} \Spec(S)^{\circ} \simeq \bigcup_{f \in \bigoplus_{n>0} S_n} U(\phi(f)).
\]
注意 $U(\phi)$ 含于 $\Spec(S')^{\circ}$。此外，它是
\[
U_+(\phi) := \bigcup_{f \in \bigoplus_{n>0} S_n} U_+(\phi(f)) \subseteq \Proj(S')
\]
在典范态射 $\Spec(S')^{\circ} \to \Proj(S')$ 下的原像。
\end{notation}

下述结果由构造 \ref{con:25.2.2} 得到。

\begin{proposition}\label{prop:25.3.2}
存在唯一的虚线态射使下述图表交换：
\[
\begin{tikzcd}[column sep=1.6em, row sep=2.6em]
\Spec(S')^{\circ} & U(\phi) \arrow[l, hook] \arrow[r] \arrow[d] & \Spec(S)^{\circ} \arrow[d, "\pi"] \\
\Proj(S') \arrow[u, "\pi'"] & U_+(\phi) \arrow[l, hook] \arrow[r, dotted] & \Proj(S)
\end{tikzcd}
\]
\end{proposition}

\begin{exercise}\label{exc:25.3.3}
设 $S$ 为 $\Z_{\ge 0}$-分次环，$d$ 为正整数。设 $S' \subseteq S$ 为由全部元素 $f \in S_{nd}$（$n \ge 0$）生成的子环。证明 $j \colon S' \to S$ 诱导同构 $\Proj(S) \simeq U_+(j) \simeq \Proj(S')$。
\end{exercise}

\begin{exercise}\label{exc:25.3.4}
设 $\phi \colon S \to S'$ 是分次同态，且当 $n$ 充分大时 $\phi_n \colon S_n \to S'_n$ 是同构。证明它诱导同构 $\Proj(S) \simeq U_+(\phi) \simeq \Proj(S')$。
\end{exercise}

\begin{lemma}\label{lem:25.3.5}
设 $A \to B$ 是 $\CRing$ 中的同态，$S$ 是 $\Z_{\ge 0}$-分次 $A$-代数。注意 $S \otimes_A B$ 自然地是 $\Z_{\ge 0}$-分次 $B$-代数。考虑 $\Z_{\ge 0}$-分次同态 $\phi \colon S \to S \otimes_A B$。则 $U_+(\phi) = \Proj(S \otimes_A B)$，且下述图表是笛卡儿方块：
\[
\begin{tikzcd}[column sep=3em, row sep=2.6em]
\Proj(S \otimes_A B) \arrow[r] \arrow[d] & \Proj(S) \arrow[d] \\
\Spec(B) \arrow[r] & \Spec(A).
\end{tikzcd}
\]
\end{lemma}

\begin{proof}
第一个断言来自：理想 $(S \otimes_A B)_+$ 由 $\phi(S_+)$ 生成。

对 $\Proj(S)$ 的任意标准开子概形 $U_+(f)$，我们有笛卡儿方块
\[
\begin{tikzcd}[column sep=3em, row sep=2.6em]
U_+(f \otimes 1) \arrow[r] \arrow[d] & U_+(f) \arrow[d] \\
\Proj(S \otimes_A B) \arrow[r] & \Proj(S).
\end{tikzcd}
\]
\end{proof}

因此只需证明
\[
\begin{tikzcd}[column sep=3em, row sep=2.6em]
U_+(f \otimes 1) \arrow[r] \arrow[d] & U_+(f) \arrow[d] \\
\Spec(B) \arrow[r] & \Spec(A)
\end{tikzcd}
\]
是笛卡儿的。这确实成立，因为
\[
U_+(f \otimes 1) \simeq \Spec\bigl((S \otimes_A B)_{(f \otimes 1)}\bigr) \simeq \Spec\bigl(S_{(f)} \otimes_A B\bigr) \simeq U_+(f) \times_{\Spec(A)} \Spec(B).
\]
\eop

\subsection{$\Proj(S)$ 的性质}\label{sec:25.4}

本小节固定一个 $\Z_{\ge 0}$-分次环 $S$。

下述结果都可通过检查标准开子集得到，我们留给读者。

\begin{lemma}\label{lem:25.4.1}
设 $S_+$ 不含幂零元。则：
\begin{enumerate}[label=(\arabic*),leftmargin=2.4em,itemsep=0.25ex]
  \item 概形 $\Proj(S)$ 既约。
  \item 概形 $\Proj(S)$ 整，当且仅当 $S_+$ 不含零因子。
\end{enumerate}
\end{lemma}

\begin{lemma}\label{lem:25.4.2}
概形 $\Proj(S)$ 拟紧，当且仅当存在有限生成的齐次理想 $I \subseteq S_+$ 使 $S_+ \subseteq \sqrt{I}$。
\end{lemma}

\begin{lemma}\label{lem:25.4.3}
概形 $\Proj(S)$ 分离。
\end{lemma}

\begin{lemma}\label{lem:25.4.4}
若 $S$ 是诺特环，则 $\Proj(S)$ 是诺特概形。
\end{lemma}

\begin{remark}\label{rem:25.4.5}
已知：$S$ 是诺特环，当且仅当 $S_0$ 是诺特环且 $S_+$ 作为 $S$ 的理想有限生成。
\end{remark}

\begin{lemma}\label{lem:25.4.6}
若 $S$ 在 $S_0$ 上有限生成，则态射 $\Proj(S) \to \Spec(S_0)$ 是有限型的。
\end{lemma}

\begin{yizhu}
引理 \ref{lem:25.4.3} 的证明并非同义反复：分离性等价于对角线 $\Proj(S) \to \Proj(S) \times_\Z \Proj(S)$ 是闭浸入，而后者只需在一对标准开集 $D_+(f)$、$D_+(g)$ 上检验。此时 $D_+(f) \cap D_+(g) = D_+(fg)$ 仿射，剩下要证 $S_{(f)} \otimes_\Z S_{(g)} \to S_{(fg)}$ 满，这由把 $h/(f^n g^m)$ 配成公分母得到。
分离性也可用来组织粘接：只要每一片仿射开子概形都分离，沿开子集粘接所得的概形就分离。

[来源：Stacks Project, Tag 01MC（Lemma 27.8.8：$\Proj(S)$ 分离及其证明），\url{https://stacks.math.columbia.edu/tag/01MC}（2026-11-17 访问）。]
\end{yizhu}

\begin{yizhu}
引理 \ref{lem:25.4.1} 的第一条值得与定义 \ref{def:25.1.2} 对照：条件只加在 $S_+$ 上，故 $S_0$ 若有幂零元，$\Proj(S)$ 仍既约。也就是说 $\Proj$ 构造完全看不见 $S_0$ 的幂零结构；这与仿射情形 $\Spec(S_0)$ 恰好相反，也是警告 \ref{warn:25.2.9} 的一个具体表现。
[依据：原书 p.123 引理 25.4.1 与 p.121 警告 25.2.9 互见。]
\end{yizhu}
